\lecture{1}
\title{Introduction}

\begin{document}

\subsection{Logistics}

The assignments for 6.1810 are:
\begin{itemize}
    \item Labs (65\%) roughly once every week, with checkoffs.
    \item A homework reading (5\%) due before the start of each lecture.
    \item One midterm and one final (30\%).
\end{itemize}

\subsection{Operating Systems}

An operating system owns:
\begin{itemize}
    \item A file system (memory of files).
    \item A network stack (communicates with the internet).
    \item Process management.
    \item State for itself and for the programs it manages.
\end{itemize}
The kernel is the ``computer program'' that implements the OS
and bridges applications in the user space with the hardware.
Applications in the user space interact with the kernel
via \emph{system calls}.

We'll be studying the xv6 operating system,
built around the RISC-V instruction set.
It's very small; all of its source code
will become completely readable by the end of this class.
This has to be run in a computer emulator (QEMU).

\subsection{File Read and Write}

\begin{verbatim}[c]
// ex1.c: copy input to output.

#include "kernel/types.h"
#include "user/user.h"

int
main()
{
  char buf[64];

  while(1){
    int n = read(0, buf, sizeof(buf));
    if(n <= 0)
      break;
    write(1, buf, n);
  }

  exit(0);
}
\end{verbatim}

Here, \verb[c]|read| and \verb[c]|write| are system calls to the kernel.
These take arguments, just like any other function call.
\begin{itemize}
    \item \verb[c]|read(fd, buf, size)|
    reads up to \verb[c]|size| bytes
    from the file with \emph{file descriptor}
    \verb[c]|fd| into the buffer at \verb[c]|buf|.
    It returns the number of bytes actually read.
    \item \verb[c]|write(fd, buf, size)|
    writes \verb[c]|size| bytes from \verb[c]|buf|
    to the file with file descriptor \verb[c]|fd|.
    \item The kernel contains, for each process,
    a table that maps file descriptors to files.
    A single file descriptor can refer to different files
    depending on which process makes the system call.
    \begin{itemize}
        \item File descriptor \verb[c]|0| is standard input.
        \item File descriptor \verb[c]|1| is standard output.
    \end{itemize}
\end{itemize}
Running this program
waits for the user to enter some input,
which the program then echoes back.

\subsection{File Creation}

\begin{verbatim}[c]
// ex2.c: create a file, write to it.

#include "kernel/types.h"
#include "user/user.h"
#include "kernel/fcntl.h"

int
main()
{
  int fd = open("out", O_WRONLY | O_CREATE | O_TRUNC);

  printf("open returned fd %d\n", fd);

  write(fd, "ooo\n", 4);

  exit(0);
}
\end{verbatim}

The \verb[c]|open(filename, flags)| system call
opens the file with name \verb[c]|filename|
and returns its file descriptor.
The second parameter \verb[c]|flags| modifies
the behavior of \verb[c]|open| like so:
\begin{itemize}
    \item The flag \verb[c]|O_CREATE| indicates that
    \verb[c]|open| should create the file
    if it does not already exist.
    \item The flag \verb[c]|O_WRONLY| indicates that
    the file should be opened for writing only.
    \item The flag \verb[c]|O_TRUNC| indicates that
    if the file already exists,
    its contents should be cut to zero bytes first.
\end{itemize}

Running this program will create a new file named
\verb[c]|out| that contains the text \verb[c]|ooo|.

\begin{remark}
    Every system call can raise an error,
    so there should always be some code
    that checks the return value of any system call
    and prints an error message accordingly
    (e.g., if \verb[c]|fd < 0|).
\end{remark}

\subsection{Process Creation}

\begin{verbatim}[c]
// ex3.c: create a new process with fork()

#include "kernel/types.h"
#include "user/user.h"

int
main()
{
  int pid;

  printf("before\n");
  pid = fork();

  printf("fork() returned %d\n", pid);

  if(pid == 0){
    printf("child\n");
  } else {
    printf("parent\n");
  }
  
  exit(0);
}
\end{verbatim}

The \verb[c]|fork| system call
creates a \emph{copy} of the running process
(in this case, \verb[c]|ex3|).
\begin{itemize}
    \item The original copy is called the \emph{parent}.
    \begin{itemize}
        \item The parent has some process ID;
        in the code above, that PID is never read.
    \end{itemize}
    \item The new copy is called the \emph{child}.
    \begin{itemize}
        \item The child has some other process ID;
        let's say it's \verb[c]|24|.
    \end{itemize}
    \item The \verb[c]|fork| system call
    returns a different thing
    depending on which process is calling it.
    \begin{itemize}
        \item In the parent process,
        \verb[c]|fork| returns the child's process ID.
        \item In the child process,
        \verb[c]|fork| will always return \verb[c]|0|.
    \end{itemize}
    It's designed this way so that a process
    can easily tell whether it's the child.
\end{itemize}
Running this program will print something like this:
\begin{verbatim}
before
fofrokk(())  rreettuurrnneedd  024

chPiAlRdE
NT
\end{verbatim}
The two processes' output interleaves character by character.
Notice how \verb[c]|before| only appears once,
because only the parent process runs that line.

\begin{remark}
    The copied process inherits its parent's memory,
    but they do not \emph{share} memories.
\end{remark}

\subsection{Replacing Processes with Executable Files}

\begin{verbatim}[c]
// ex4.c: replace a process with an executable file

#include "kernel/types.h"
#include "user/user.h"

int
main()
{
  char *argv[] = { "echo", "this", "is", "echo", 0 };

  exec("echo", argv);

  printf("exec failed!\n");

  exit(0);
}
\end{verbatim}

Here, \verb[c]|echo| is a program pre-compiled into xv6
(the executable \verb[c]|/echo|)
that prints its arguments back out,
separated by spaces, and then a newline.

The \verb[c]|exec(filename, args)| system call
tells the kernel to completely discard the calling process's memory
and instead run the program \verb[c]|filename|;
if \verb[c]|exec| succeeds, nothing after it
in the calling program ever runs.

Naturally, \verb[c]|args| is the list of arguments
to the program \verb[c]|filename|.
The first element in \verb[c]|args| is not printed by \verb[c]|echo|;
by convention, the first element is always the name of the program
to be run.

Running \verb[c]|ex4| prints the following.

\begin{verbatim}
this is echo

\end{verbatim}

\subsection{Waiting}

\begin{verbatim}[c]
// ex5.c: fork then exec

#include "kernel/types.h"
#include "user/user.h"

int
main()
{
  int pid, status;

  pid = fork();
  if(pid == 0){
    char *argv[] = { "echo", "THIS", "IS", "ECHO", 0 };
    exec("echo", argv);
    printf("exec failed!\n");
    exit(1); // this line never gets called
  } else {
    printf("parent waiting\n");
    int xxx = wait(&status);
    printf("child %d exited with status %d\n", xxx, status);
  }

  exit(0);
}
\end{verbatim}

The \verb[c]|wait(&status)| system call
makes a parent process pause until
one of its children exits,
at which point \verb[c]|wait| will:
\begin{itemize}
    \item Return the PID \verb[c]|xxx|
    of the child that exited.
    \item Update the data
    at address \verb[c]|&status|
    with the way the child exited.
    \begin{itemize}
        \item If the child exits via \verb[c]|exit(3);|,
        then \verb[c]|3| is written into \verb[c]|status|.
    \end{itemize}
\end{itemize}

Running this program produces the following.

\begin{verbatim}
parent waiting
THIS IS ECHO
child 11 exited with status 0
\end{verbatim}

Notably, the \verb[c]|echo| process
exits via \verb[c]|exit(0);|.

\subsection{Output Redirecting}

\begin{verbatim}[c]
// ex6.c: run a command with output redirected

#include "kernel/types.h"
#include "user/user.h"
#include "kernel/fcntl.h"

int
main()
{
  int pid;

  pid = fork();
  if(pid == 0){
    close(1);
    open("out", O_WRONLY | O_CREATE | O_TRUNC);

    char *argv[] = { "echo", "ex6's", "redirected", "echo", 0 };
    exec("echo", argv);
    printf("exec failed!\n");
    exit(1);
  } else {
    wait((int *) 0); // this is a null pointer: "don't write the status"
    printf("wait returned\n");
  }

  exit(0);
}
\end{verbatim}

The system call \verb[c]|close(fd);| just
closes the file with file descriptor \verb[c]|fd|.

The goal, here, is to somehow replace the standard output file
with the \verb[c]|out| file created by the child process.
In other words, the goal is to take
the child process's file-descriptor-to-file table,
delete the row that reads \verb|1 = standard-output|,
and add a row that reads \verb|1 = out|.
This way, the \verb[c]|echo| program
will write to \verb|1 = out|.

The code above relies on the following trick:
the system call \verb[c]|open(filename, flags)|
always assigns the file \verb[c]|filename|
to the \emph{lowest} available file descriptor.
Since file descriptor \verb[c]|1| was \emph{just} deleted,
it is guaranteed that the kernel
will create file \verb[c]|out| so that
its file descriptor is \verb[c]|1|.

\begin{remark}
    Note that this reassignment of
    file descriptor \verb[c]|1| does \emph{not} affect
    the parent process's standard output.
    This is because---repeating this again---the kernel
    keeps separate file-descriptor-to-file tables
    for each individual process.
\end{remark}

Here's what calling this process
in the terminal might look like.

\begin{verbatim}
$ ex6
wait returned
$ cat out
ex6's redirected echo
\end{verbatim}

This effectively implements IO redirection,
which can also be accomplished in the shell
using the \verb|>| symbol.

\begin{verbatim}
$ echo ex6's redirected echo > out
$ cat out
ex6's redirected echo
\end{verbatim}

Only the shell needs to know how to interpret \verb|>|.

\end{document}